\section{Introduction}

Per- and polyfluoroalkyl substances (PFAS) are persistent, mobile synthetic chemicals now regulated in
drinking water at the nanogram-per-litre level \citep{buck2011, guelfo2021, epa2024mcl}. The PFAS
measured in an in-stream grab sample is the sum of several transport pathways acting at once --- surface
runoff, lateral and sediment-bound transport, biosolids and industrial inputs, and the discharge of
contaminated groundwater \citep{prevedouros2006, sepulvado2011}. Because PFAS management now turns on
\emph{where} the in-stream contamination originates --- a Clean Water Act load allocation, a CERCLA cost
recovery, a source-control or remediation decision \citep{epa2024cercla, egle2022rule57} --- the
governing scientific need is no longer to detect PFAS but to attribute an observed in-stream
concentration to the pathways that produced it.

Among those pathways, groundwater discharge is at once the most consequential and the least resolvable.
Where a legacy source exists, PFAS that infiltrated decades ago migrates through the aquifer and
re-enters the stream as baseflow, so that a substantial share of the in-stream load can arrive through
the streambed rather than across the land surface \citep{rafiei2023pfas}. Yet this is exactly the
pathway that observation cannot disaggregate: a grab sample cannot separate groundwater-borne from
surface-derived PFAS, and the statistical occurrence models that dominate large-scale PFAS prediction
return a probability of contamination, not a mass-conserving, pathway-resolved concentration
\citep{tokranov2024groundwater}. Quantifying the groundwater contribution to in-stream PFAS therefore
requires a process model that carries the contaminant continuously from the land surface, through the
aquifer, and back to the channel.

Watershed-scale PFAS modeling has been advancing toward this goal. PFAS was first simulated in receiving
rivers on a hydrology supplied by a coupled surface-water/groundwater model \citep{raschke2022integrated},
and a distributed watershed model subsequently represented PFAS fate and transport across the land
surface and identified groundwater as a watershed-scale source \citep{rafiei2023pfas}. These advances
established that process-based watershed PFAS modeling is tractable and that the subsurface is material.
They share one structural limit: the contaminant is passed in a single direction, from the surface
model into the subsurface, so the PFAS that a legacy plume discharges back to the stream is never
returned to the surface-water mass balance. The groundwater \emph{contribution} to the in-stream load
--- the very quantity that source attribution depends on --- thus remains outside the model.

We advance watershed-scale PFAS fate and transport by resolving this missing direction: the two-way
exchange of PFAS between surface water and groundwater. A process-based, three-phase soil model carries
PFAS from the soil profile to the stream, partitioning among the aqueous phase, the Freundlich solid
phase \citep{higgins_luthy2006}, and the air--water interface that governs retention in unsaturated
media \citep{costanza2019, brusseau2018, guo2020}; a groundwater flow-and-transport model then carries
the infiltrated PFAS through the aquifer and returns the groundwater-borne fraction to the channel,
closing a single continuous PFAS mass balance across the two compartments. The subsurface leg is built
on MODFLOW~6, whose native Freundlich groundwater transport and advanced streamflow transport
\citep{langevin2017mf6, morway2021gwt} provide the contaminant physics that the embedded-solver
couplings underlying earlier watershed models do not carry \citep{bailey2016swatmodflow,
bailey2025swatplusmodflow}; fully integrated hydrologic simulators exist outside the watershed-modeling
family but do not carry a watershed's surface PFAS budget \citep{markstrom2008gsflow, brunner2012hgs}.

This capability lets us pose a question the one-directional models could not answer: at an impacted
watershed, how much of the in-stream PFAS is groundwater-borne, and where along the river does it enter?
We address it on the Rogue River, Michigan, downstream of the Wolverine World Wide / House Street legacy
PFAS source --- a tannery-waste disposal area on a groundwater divide immediately upgradient of the
river, where source-area groundwater has reached the order of $10^{5}$~ng~L\textsuperscript{-1} and
where a groundwater interceptor was later engineered specifically to capture PFAS before it discharges
to the channel \citep{epa_wolverine, egle_housestreet, egle_rrt5_2021}. The site offers an unusually
rich public record --- thousands of water-table observations, hundreds of groundwater PFAS
measurements, streambed pore water sampled at the discharge interface, and a multi-analyte in-stream
network \citep{egle_grandriver2019, lemke2024} --- against which every compartment of the model can be
tested (study-area configuration in Table~\ref{tab:studyarea}).

Two independent lines of evidence then converge on the answer. A model-free, multi-analyte fingerprint
decomposition of the in-stream signal shows that the rising PFOS gradient along the lower mainstem
carries the chemical signature of the groundwater plume; and a joint surface/groundwater calibration ---
constructed to avoid the double-counting that a naive sum of the two legs would incur --- attributes a
first-order, spatially localized share of the lower-mainstem load to legacy groundwater discharge. We
further report the first systematic global sensitivity and uncertainty analysis of coupled watershed
PFAS fate and transport, isolating the parameters that govern how PFAS is partitioned between the
surface and subsurface routes. The result is a pathway-resolved, mass-conserving account of PFAS fate
and transport across the surface-water and groundwater compartments of a single watershed --- the basis
for attributing in-stream PFAS to its sources and for evaluating the pathway-specific actions that PFAS
management requires.
